\subsection{The remaining ordinary RH Selmer factors}
\label{sec:rh-selmer-low-dyadic}

\begin{corollary}[All exponent-one residual factors for RH]
\label{cor:rh-selmer-exponent-one}
For every odd $\ell\nmid N$ and every $e\ge1$ in the original
twist range, the actual character bound (DC2) holds whenever
$2^b>8\ell^2$, where $2^b$ is its dyadic conductor.
In particular the complete ordinary RH large-conductor
Case~V estimate at three is proved.

For the RH residual product of
Proposition~\ref{shared-rh-integral-product}, there is an
absolute constant $B$ such that the entire exponent-one
contribution satisfies, for every choice of its original
weights $0\le w_\ell\le1$,
\[
       \sum_{\ell\ne2}w_\ell(\log\ell)
                       \operatorname{length}_{\mathbf Z_\ell} M_{\ell,1}
                         \le (2N)^B .                          \tag{DC10}
\]
Here the notation $M_{\ell,1}$ is that of the residual product;
only its finite supported factors are summed.
\end{corollary}

\begin{proof}
Apply Theorem~\ref{thm:rh-dyadic-host-count}.
At $e=1$ the original threshold is $9216>72$, so every
required large-conductor factor is included. Summing over at
most $N$ characters and using the exact integral Fourier
decomposition gives the bound at three for $T_{3,1}$ in (RN6).
For all $\ell\ge5$ the exponent one is nonresonant, and
Theorem~\ref{thm:rh-selmer-all-nonresonant} gives the aggregate
bound, including primes in the exceptional set and all of
their character factors. Multiplying each nonnegative summand
by its weight cannot increase the sum. Enlarge the exponent
to absorb the fixed constants. This proves (DC10).
\end{proof}

\begin{corollary}[Reduction of the higher resonant factors]
\label{cor:rh-selmer-low-dyadic}
For $e\ge2$ and $\ell-1\mid2e$ with $\ell$ odd, put
\[
 c_\ell=\lfloor\log_2(8\ell^2)\rfloor,\qquad
 E_{\ell,\mathrm{low}}=
 \mathbf Q(\mu_{2^{\min(a,c_\ell)}})^+
          (\sqrt{-q},\sqrt{-d}),\qquad
 H_\ell=\operatorname{Gal}(E/E_{\ell,\mathrm{low}}).
\]
Only the character factors of
\[
 C_e(1-A_\Delta)A_{H_\ell}S_{\ell,e}                 \tag{DC11}
\]
remain in (RN5). The omitted factors, summed over all
resonant primes for a fixed $e$, have total logarithmic
order bounded by $[N(e+1)]^B$ for an absolute exponent $B$,
also with their original weights.
Moreover
\[
 [E_{\ell,\mathrm{low}}:\mathbf Q]\le8\ell^2
                             \le8(2e+1)^2 .
\]
The invariant factors identify with the corresponding
strict groups of the reduced fields by restriction.
The two varying quadratic factors are retained.
\end{corollary}

\begin{proof}
The size conditions of (RN4) give $\ell\nmid N$.
In the integral Fourier decomposition, available because
the Galois group is a $2$-group and $\ell$ is odd, a
character has dyadic conductor at most $2^{c_\ell}$ if
and only if it is trivial on $H_\ell$. The invariant
part is the image of $A_{H_\ell}$.
Every other character has conductor $2^b>8\ell^2$,
so (DC2) applies to its original strict group.
Use the normalization and coefficient-field comparison
of (RN2)--(RN5). There are at most $N$ characters and
at most $2e+1$ resonant primes, each at most $2e+1$;
summing gives the asserted bound. The weights lie in
$[0,1]$ and cannot increase it. The field degree is
at most four times the real dyadic cyclotomic degree,
hence at most $2^{c_\ell}\le8\ell^2$.

The extension $E/E_{\ell,\mathrm{low}}$ has degree prime
to $\ell$. Inflation-restriction identifies cohomology
of the reduced field with the $H_\ell$-invariants for
the original discrete $\ell$-primary coefficients.
At every place restriction on the local cohomology is
injective since the local index is prime to $\ell$.
The strict local kernels are therefore preserved in
both directions. The parity and trivial-character
projectors commute with restriction, proving (DC11).
\end{proof}

The next obligation is a uniform integral estimate for
(DC11), not just a bound for the field degree, followed
by the original weighted assembly. In particular the
remaining dyadic conductors at three are at most $64$.
No general Dirichlet conductor reduction is inferred
from this RH-specific decomposition.
